\documentclass[11pt, a4paper, oneside]{report}
\usepackage[UTF8]{ctex}
\usepackage{graphicx}
\usepackage{geometry}
\usepackage{titlesec}
\usepackage{fancyhdr}
\usepackage{listings}
\usepackage{xcolor}
\usepackage{amsmath, amssymb}
\usepackage{hyperref}
\usepackage{booktabs}
\usepackage{array}
\usepackage{longtable}
\usepackage{float}
\usepackage{caption}
\usepackage{subcaption}
\usepackage{url}

% ==========================================
% 页面布局与基本设置
% ==========================================
\geometry{margin=2.5cm}
\renewcommand{\baselinestretch}{1.2}

\setcounter{tocdepth}{3}
\setcounter{secnumdepth}{3}

% 章节标题样式配置
\titleformat{\chapter}[display]{\normalfont\huge\bfseries}{\chaptertitlename\ \thechapter}{20pt}{\Huge}
\titleformat{\section}[hang]{\normalfont\Large\bfseries}{\thesection}{1em}{}
\titleformat{\subsection}[hang]{\normalfont\large\bfseries}{\thesubsection}{1em}{}

% ==========================================
% 页眉页脚配置
% ==========================================
\pagestyle{fancy}
\fancyhf{}
\chead{\leftmark}
\cfoot{\thepage}
\renewcommand{\headrulewidth}{0.4pt}
\renewcommand{\footrulewidth}{0.4pt}
\fancypagestyle{plain}{
    \fancyhf{}
    \cfoot{\thepage}
    \renewcommand{\headrulewidth}{0pt}
}

% ==========================================
% 代码清单样式（GLSL / JavaScript）
% ==========================================
\definecolor{codegray}{rgb}{0.5,0.5,0.5}
\definecolor{codegreen}{rgb}{0.0,0.5,0.0}
\definecolor{codeblue}{rgb}{0.1,0.2,0.7}
\definecolor{codebg}{rgb}{0.97,0.97,0.97}
\lstdefinestyle{shaderstyle}{
    backgroundcolor=\color{codebg},
    commentstyle=\color{codegreen},
    keywordstyle=\color{codeblue}\bfseries,
    numberstyle=\tiny\color{codegray},
    basicstyle=\ttfamily\footnotesize,
    breaklines=true,
    captionpos=b,
    numbers=left,
    numbersep=6pt,
    showstringspaces=false,
    tabsize=2,
    frame=single,
    rulecolor=\color{codegray}
}
\lstset{style=shaderstyle}

% 自定义快捷命令
\newcommand{\keyword}[1]{\textbf{#1}}
\newcommand{\code}[1]{\texttt{#1}}
% 图片守卫：存在则插图，缺失则显示占位框（便于无截图时也能编译）
\newcommand{\figbox}[1]{%
  \IfFileExists{#1}{\includegraphics[width=\textwidth]{#1}}%
  {\fbox{\parbox[c][3.6cm][c]{0.88\textwidth}{\centering 占位图，请将截图保存为 \texttt{#1}}}}%
}

% ==========================================
% 文档基本信息
% ==========================================
\title{计算机图形学与高级渲染期末大作业报告}
\author{宋佳刚}
\date{\today}

\begin{document}

% ------------------------------------------
% 封面页
% ------------------------------------------
\begin{titlepage}
    \centering
    \vspace*{2cm}
    {\Huge\bfseries 计算机图形学与高级渲染期末大作业\par}
    \vspace{1.5cm}
    {\LARGE 基于 WebGL/Three.js 的室内交互式高级渲染系统：\\[0.4em]Blinn-Phong 光照、法线贴图与阴影映射的原理与实现\par}

    \vspace{3cm}
    {\large
     \begin{tabular}{ll}
        \textbf{专业：} & \underline{\hspace{4em}数字媒体技术\hspace{4em}} \\[0.8em]
        \textbf{学号：} & \underline{\ 240207132\ } \\[0.8em]
        \textbf{姓名\&工作占比：} & \underline{\ 宋佳刚 \quad 50\%\ } \\[0.8em]
        \textbf{同组同学\&工作占比：} & \underline{\ 赵志国 \quad 50\%\ }\\[0.8em]
        \textbf{完成日期：} & \today \\[0.8em]
    \end{tabular}
    \par}
    \vfill
\end{titlepage}

% ------------------------------------------
% 摘要页
% ------------------------------------------
\chapter*{摘要}
\addcontentsline{toc}{chapter}{摘要}
本项目遵循“理论指导实践”的图形学研究方法，构建了一个运行于浏览器、基于 WebGL 与 Three.js 的室内交互式三维场景，系统实现并验证了经典局部光照模型、表面细节增强与实时阴影等关键渲染技术。区别于直接调用引擎黑盒材质的做法，本文从数学原理出发，重新推导并以原生 GLSL 着色器白盒化实现了 Blinn-Phong 光照模型，使其可与基于物理的渲染（PBR）一键切换对比；同时通过切线空间法线贴图在不增加几何面数的前提下表现表面微观凹凸，并以两遍式阴影映射（Shadow Mapping）结合百分比邻近过滤（PCF）生成柔和阴影。

在几何资产层面，场景中全部三维模型均在 Blender 中以 Python 脚本参数化建模，导出 glTF 二进制（GLB）后由 GLTFLoader 载入，几何类型覆盖平面、立方体、球体、圆柱、圆锥及自定义合并网格。系统进一步实现了第一人称漫游、透视/正交投影切换、屏幕射线拾取与高亮、实时参数面板，以及基于帧缓冲（FBO）的 Bloom 辉光与 Sobel 边缘检测两类后处理特效。实验表明，本系统在保持实时帧率（约 60--165 FPS）的同时，能清晰呈现光照、阴影与表面细节，验证了从数学模型到实时渲染的高效转化。本项目是对课程两次实验所学的综合运用：在沿用实验一光栅化管线（坐标变换、纹理映射与抗锯齿）的基础上，借鉴实验二光线追踪与蒙特卡洛全局光照中的 BRDF 与可见性理论，对全局光照效果进行实时近似。

\textbf{关键词：} WebGL, Three.js, Blinn-Phong 光照, 法线贴图, 阴影映射(Shadow Mapping), PCF 软阴影, 后处理(FBO), Blender 程序化建模

\tableofcontents

% ------------------------------------------
% 第一章：绪论
% ------------------------------------------
\chapter{绪论}

\section{课题研究背景与意义}
随着实时渲染技术在游戏、虚拟现实与数字媒体领域的广泛应用，如何在有限的计算预算下实现兼具真实感与交互性的三维场景，成为计算机图形学的核心议题。现代图形管线以可编程着色器为核心，开发者需要深入理解光照模型、纹理映射与阴影算法的数学原理，才能在效果与性能之间取得平衡。本课题以一个室内场景为载体，强调“理论先行、白盒实现”，不满足于直接套用引擎内置材质，而是从公式推导出发，亲手实现 Blinn-Phong 光照、法线贴图与阴影映射等关键技术，以加深对实时渲染管线的理解。

\section{国内外研究现状与主流技术对比}
为明确本项目的技术定位，本节从“光照着色模型”与“几何与阴影管线”两个维度，将本项目所采用的技术与主流方案进行横向对比。

\subsection{经验光照模型与基于物理的渲染(PBR)的差异讨论}
\begin{enumerate}
    \item \textbf{经验模型（Phong / Blinn-Phong）}：以 Lambert 漫反射加经验性镜面高光项描述表面反射，参数（漫反射系数、镜面系数、高光指数）依赖手工调节，计算量小、实时性好，但不满足能量守恒，在不同光照条件下一致性较差。
    \item \textbf{基于物理的渲染（PBR）}：以微表面理论与 Cook-Torrance BRDF 为基础，使用金属度（metalness）、粗糙度（roughness）等具有物理含义的参数，满足能量守恒，跨光照条件表现稳定，已成为现代引擎的事实标准，但单次着色开销更高。
\end{enumerate}
本项目同时实现了二者：以 Three.js 的 \code{MeshStandardMaterial} 提供 PBR 通路，并自写 GLSL Blinn-Phong 着色器作为对照，支持运行时一键切换，直观展示经典经验模型与现代物理模型的视觉差异（见第三、五章）。

\subsection{光栅化阴影与光线追踪阴影的差异讨论}
\begin{enumerate}
    \item \textbf{Shadow Mapping（光栅化阴影）}：以光源视角渲染深度图，再在相机视角下做深度比较判断遮挡，复杂度低、实时性强，是当前实时渲染的主流；其缺陷（走样、自遮挡）需借助 PCF、偏移（bias）等手段缓解。
    \item \textbf{Ray-Traced Shadow（光线追踪阴影）}：向光源投射阴影射线求交，物理正确、软硬阴影自然，但开销大，依赖硬件光追加速。
\end{enumerate}
综合实时性与课程要求，本项目采用 Shadow Mapping 配合 PCF 软阴影（见第四章），并通过参数面板提供阴影开关以对比其渲染开销。

\subsection{实时光栅化管线与蒙特卡洛路径追踪的差异讨论}
本课程的两次实验恰好覆盖了实时渲染的两条主线，可作为本项目技术定位的重要参照：\keyword{实验一}围绕\keyword{光栅化（Rasterization）管线}展开，依次实现了模型变换与齐次坐标（MVP）、三角形光栅化与重心坐标插值、超采样抗锯齿（SSAA）、纹理映射与 Mipmap 纹理抗锯齿；\keyword{实验二}则进入\keyword{光线追踪与蒙特卡洛全局光照}，实现了 Möller--Trumbore 三角形求交、BVH 加速结构、漫反射 BSDF（$f_r=\rho/\pi$）、直接光照的半球均匀采样与光源重要性采样，以及带下一事件估计（NEE）与俄罗斯轮盘赌（Russian Roulette）的路径追踪。二者的本质差异在于：
\begin{enumerate}
    \item \textbf{光栅化（本项目所属范式）}：以“图元投影 + 深度测试 + 片元着色”为核心，逐三角形并行处理，配合 Shadow Mapping、环境光、Bloom 等手段\keyword{近似}全局光照效果，复杂度低、可实时交互，是游戏与可视化的主流。
    \item \textbf{蒙特卡洛路径追踪（课程实验二）}：以光线-场景求交与对渲染方程的随机采样积分为核心，物理正确地求解多重反弹、软阴影与间接光照，真实度高但收敛慢、单帧开销大，多用于离线渲染。
\end{enumerate}
本项目立足于\keyword{实时光栅化}一极，但在报告与实现中自觉借鉴实验二的理论：例如环境光项 $k_a I_a$ 是对间接光照（全局光照）的廉价近似，PCF 软阴影是对路径追踪中可见性项 $V(p,p')$ 软过渡的经验模拟，从而在实时帧率下逼近离线渲染的视觉效果。

\subsection{适用场景横向对比}
如表~\ref{tab:compare} 所示，不同技术路线各有侧重。

\begin{table}[htbp]
\centering
\caption{本项目所用技术与主流方案的适用性对比}
\label{tab:compare}
\begin{tabular}{p{3cm}|p{5.5cm}|p{5.5cm}}
\hline
\textbf{维度} & \textbf{本项目方案} & \textbf{对照/工业主流} \\ \hline
\textbf{光照模型} & 自写 Blinn-Phong（白盒、可教学）+ PBR 对照 & 引擎黑盒 PBR（Cook-Torrance） \\ \hline
\textbf{阴影} & Shadow Mapping + PCF 软阴影 & 光线追踪阴影 / VSM / CSM \\ \hline
\textbf{表面细节} & 程序化切线空间法线贴图 & 美术烘焙法线/位移贴图 \\ \hline
\textbf{几何资产} & Blender 脚本程序化建模 + GLB 加载 & DCC 手工建模 + 流水线 \\ \hline
\end{tabular}
\end{table}

\section{本项目与课程实验内容的关联}
本项目并非孤立选题，而是对课程两次实验所学知识的综合运用与延伸。表~\ref{tab:course} 将本项目的核心技术逐项映射到对应的课程实验内容，体现“结合课程所学”的设计思路。

\begin{table}[htbp]
\centering
\caption{本项目核心技术与课程实验内容的对应关系}
\label{tab:course}
\begin{tabular}{p{4cm}|p{4.2cm}|p{4.2cm}}
\hline
\textbf{本项目技术} & \textbf{对应课程实验内容} & \textbf{关联说明} \\ \hline
MVP 模型-视图-投影变换、透视/正交切换 & 实验一：坐标变换、模型变换（齐次坐标与矩阵乘法） & 同一套齐次坐标与矩阵级联理论，本项目用于实时相机管线 \\ \hline
多纹理映射与法线贴图 & 实验一：纹理映射、Mipmap 纹理抗锯齿、纹理采样 & 沿用 UV 采样与 Mipmap 思想，并扩展到切线空间法线贴图 \\ \hline
MSAA 抗锯齿、PCF 软阴影 & 实验一：超采样抗锯齿（SSAA） & 同源的“多采样取平均”抗锯齿思想，分别用于边缘与阴影 \\ \hline
Blinn-Phong / PBR 漫反射项 & 实验二：漫反射 BSDF（$f_r=\rho/\pi$）、PBR 讲义 & 共享 Lambert 漫反射与 BRDF 理论基础 \\ \hline
环境光近似、Bloom、PCF 软过渡 & 实验二：蒙特卡洛全局光照、路径追踪 & 以经验/后处理手段实时近似离线全局光照效果 \\ \hline
三角形网格求交（射线拾取） & 实验二：Möller--Trumbore 三角形求交、BVH & 射线-三角形相交理论的实时交互应用 \\ \hline
\end{tabular}
\end{table}

\section{本课题研究内容与框架体系}
本报告后续章节编排如下：第二章阐述场景几何的 Blender 程序化建模管线与 MVP 变换；第三章推导 Blinn-Phong 光照模型与切线空间法线贴图的数学原理及 GLSL 实现；第四章阐述 Shadow Mapping 阴影映射与基于 FBO 的后处理特效；第五章介绍交互系统与实时性能量化；第六章总结全文并展望后续优化方向。

% ------------------------------------------
% 第二章
% ------------------------------------------
\chapter{场景几何建模与 MVP 变换管线}

\section{Blender 程序化建模与资产管线}
本项目坚持“所有模型自建”的原则，场景中房间、木桌、椅子、台灯、书本与装饰球等全部几何体均在 Blender 中通过 Python（\code{bpy}）脚本参数化生成，而非依赖外部模型下载。脚本以统一的基本体构造函数（立方体、圆柱、圆锥、球、平面）按尺寸与位置实例化物体，并用 \code{join} 操作将多基本体合并为“自定义网格模型”（如由桌面与四腿合并的木桌）。建模完成后，经 glTF 2.0 导出器以启用 Y-up 转换的方式导出为二进制 GLB，再由 Three.js 的 \code{GLTFLoader} 异步载入。该管线实现了“参数化建模 $\rightarrow$ 标准格式导出 $\rightarrow$ 实时引擎加载”的完整闭环，几何类型完整覆盖了课程要求的平面、立方体、球体、圆柱体与自定义网格模型。

\section{模型-视图-投影(MVP)变换}
三维顶点从模型空间到裁剪空间的变换由模型矩阵 $M$、视图矩阵 $V$ 与投影矩阵 $P$ 级联完成。对模型空间齐次坐标 $p_{\text{model}}$，其裁剪空间坐标为：
\begin{equation}
p_{\text{clip}} = P \cdot V \cdot M \cdot p_{\text{model}}
\end{equation}
其中 $M$ 由物体的平移、旋转、缩放复合而成，$V$ 为相机世界变换的逆，$P$ 决定投影方式。本项目实现了\keyword{透视投影}与\keyword{正交投影}双相机切换：两台相机共享位置与朝向，仅投影矩阵不同。透视投影矩阵将视锥映射为规则观察体并引入近大远小的透视除法；正交投影矩阵则为平行投影，保持物体尺寸与深度无关。Three.js 在渲染时自动完成上述 MVP 级联与法线矩阵的计算，开发者在自写着色器中可直接使用引擎注入的 \code{modelMatrix}、\code{viewMatrix}、\code{projectionMatrix} 等内建一致变量。该 MVP 变换理论正是课程实验一“坐标变换”与“模型变换”所学内容的直接应用——实验一中通过齐次坐标与矩阵乘法对二维/三维顶点做批量变换，本项目则将同一套理论用于实时相机管线，并扩展出透视与正交两种投影矩阵的切换。

% ------------------------------------------
% 第三章
% ------------------------------------------
\chapter{Blinn-Phong 光照模型与法线贴图的原理与实现}

\section{Blinn-Phong 局部光照模型}
Blinn-Phong 模型由 Phong 模型改进而来，将表面出射光照分解为环境光、漫反射与镜面反射三项之和。对场景中每个光源 $i$，着色点的出射光强为：
\begin{equation}
\label{eq:blinnphong}
I = k_a I_a + \sum_{i} \Big[\, k_d I_i \max(0,\, \vec{N}\cdot\vec{L}_i) + k_s I_i \max(0,\, \vec{N}\cdot\vec{H}_i)^{\alpha} \,\Big]
\end{equation}
其中 $\vec{N}$ 为表面单位法向量，$\vec{L}_i$ 为指向第 $i$ 个光源的单位向量，$\vec{V}$ 为指向观察者的单位向量，$k_a, k_d, k_s$ 分别为环境、漫反射、镜面反射系数，$\alpha$ 为高光指数，$I_a, I_i$ 为光强。环境项 $k_a I_a$ 近似全局间接光照；漫反射项遵循 Lambert 余弦定律，与 $\vec{N}\cdot\vec{L}_i$ 成正比。值得注意的是，此处的漫反射与课程实验二中的漫反射 BSDF 同源：实验二以蒙特卡洛积分求解渲染方程，其漫反射项 $f_r=\rho/\pi$ 同样建立在 Lambert 反射假设之上，本项目则将其用于实时局部光照的确定性求值。

与原始 Phong 使用反射向量 $\vec{R}$ 不同，Blinn 引入\keyword{半程向量（half-vector）} $\vec{H}$ 计算高光，避免逐像素求反射向量，且在掠射角下高光形状更自然：
\begin{equation}
\vec{H}_i = \frac{\vec{L}_i + \vec{V}}{\lVert \vec{L}_i + \vec{V} \rVert}
\end{equation}
对点光源，本项目引入平方反比衰减以模拟光强随距离衰减：
\begin{equation}
f_{\text{att}} = \frac{1}{1 + d^2 / r^2}
\end{equation}
其中 $d$ 为着色点到点光源距离，$r$ 为可调衰减半径。场景中配置\keyword{方向光（主光）}与\keyword{点光（台灯）}两个光源，满足“至少两个普通光”的要求。

\subsection{GLSL 白盒实现}
本人以 Three.js 的 \code{ShaderMaterial} 实现上述模型。顶点着色器（代码清单~\ref{lst:blinnvs}）将顶点经模型矩阵变换至世界空间，并把世界坐标与法向量通过 \code{varying} 传递给片元着色器；片元着色器（代码清单~\ref{lst:blinn}）再对方向光与点光分别按式~(\ref{eq:blinnphong}) 累加 Blinn-Phong 贡献。这里复用了第二章的 MVP 变换：\code{modelMatrix}、\code{viewMatrix}、\code{projectionMatrix} 均为 Three.js 注入的内建一致变量。

\begin{lstlisting}[language=C, caption={Blinn-Phong 顶点着色器（输出世界空间坐标与法向量）}, label={lst:blinnvs}]
varying vec3 vWorldPos;
varying vec3 vWorldNormal;
void main() {
    vec4 wp = modelMatrix * vec4(position, 1.0);        // 模型空间 -> 世界空间
    vWorldPos    = wp.xyz;
    vWorldNormal = mat3(modelMatrix) * normal;          // 法向量变换(刚体)
    gl_Position  = projectionMatrix * viewMatrix * wp;  // 世界 -> 裁剪空间(MVP)
}
\end{lstlisting}

\begin{lstlisting}[language=C, caption={Blinn-Phong 单光源贡献（GLSL 片元着色器核心）}, label={lst:blinn}]
// 单个光源的 Blinn-Phong 贡献：漫反射 + 镜面反射
vec3 lighting(vec3 N, vec3 V, vec3 L, vec3 radiance) {
    float diff = max(dot(N, L), 0.0);
    vec3  H    = normalize(L + V);
    float spec = (diff > 0.0) ? pow(max(dot(N, H), 0.0), uShininess) : 0.0;
    return radiance * (diff * uDiffuse + spec * uSpecular);
}

void main() {
    vec3 N = normalize(vWorldNormal);
    vec3 V = normalize(cameraPosition - vWorldPos);
    vec3 color = uAmbient * uDiffuse;                 // 环境项
    color += lighting(N, V, normalize(uDirDir), uDirColor);  // 方向光
    vec3  Lp   = uPointPos - vWorldPos;               // 点光方向
    float dist = length(Lp);
    float att  = 1.0 / (1.0 + (dist*dist)/(uPointDist*uPointDist));
    color += lighting(N, V, Lp/dist, uPointColor * att);     // 点光(带衰减)
    gl_FragColor = vec4(color, 1.0);  // 线性输出，交由 OutputPass 做色调映射
}
\end{lstlisting}

\section{切线空间法线贴图（Normal Mapping）}
法线贴图在不增加几何细节的前提下，通过逐像素扰动法向量表现表面微观凹凸。其核心是将以 RGB 编码法向量的纹理映射到表面，渲染时以纹理法线替代插值几何法线参与光照。法线贴图通常存储于\keyword{切线空间}，即以切线 $\vec{T}$、副切线 $\vec{B}$、法线 $\vec{N}$ 为基的局部坐标系。纹理像素范围 $[0,1]$ 需先解码为 $[-1,1]$ 的法向量：
\begin{equation}
\vec{n}_{t} = 2\cdot \text{texel}_{rgb} - 1
\end{equation}
再经 TBN 矩阵变换到世界空间参与光照：
\begin{equation}
\text{TBN} = [\,\vec{T}\ \ \vec{B}\ \ \vec{N}\,], \qquad \vec{n}_{world} = \text{normalize}(\text{TBN}\cdot \vec{n}_{t})
\end{equation}

\subsection{程序化法线贴图生成}
本项目未使用外部贴图，而是在运行时程序化生成法线贴图并应用于地面与墙面：先以随机噪声构造高度场 $h(u,v)$ 并做盒式模糊使之连续，再用 Sobel 算子求 $u,v$ 方向梯度，将梯度组合为切线空间法线后归一化写入纹理：
\begin{equation}
\vec{n}_{t} = \text{normalize}\big(\,-s\,\partial h/\partial u,\ -s\,\partial h/\partial v,\ 1\,\big)
\end{equation}
其中 $s$ 为凹凸强度系数。生成法线贴图的核心代码如代码清单~\ref{lst:normalgen} 所示：对每个纹素用中心差分近似高度场梯度 $\partial h/\partial u,\ \partial h/\partial v$，组合为切线空间法线后，按 $\,0.5\vec{n}+0.5\,$ 重映射回 $[0,1]$ 写入 RGB 通道。

\begin{lstlisting}[language=Java, caption={程序化法线贴图生成（JavaScript：高度场梯度 $\rightarrow$ 切线空间法线）}, label={lst:normalgen}]
// h[]: 经盒式模糊的连续高度场; idx(x,y): 环绕取址
for (let y = 0; y < size; y++)
  for (let x = 0; x < size; x++) {
    const gx = h[idx(x+1,y)] - h[idx(x-1,y)];   // 中心差分: dh/du
    const gy = h[idx(x,y+1)] - h[idx(x,y-1)];   // 中心差分: dh/dv
    const nx = -gx*strength, ny = -gy*strength, nz = 1.0;
    const len = Math.hypot(nx, ny, nz);          // 归一化
    const o = (y*size + x) * 4;
    img.data[o]   = (nx/len*0.5 + 0.5) * 255;    // [-1,1] -> [0,1] -> R
    img.data[o+1] = (ny/len*0.5 + 0.5) * 255;    //                    G
    img.data[o+2] = (nz/len*0.5 + 0.5) * 255;    //                    B
    img.data[o+3] = 255;
  }
\end{lstlisting}
该方法零外部资源依赖，即可在地面与墙面呈现细腻的微观起伏，在几何面数不变的情况下显著增强真实感。法线贴图的本质仍是\keyword{纹理采样}：渲染时按 UV 坐标对法线纹理采样，与课程实验一的纹理映射、Mipmap 纹理抗锯齿同属一套理论。本项目对法线纹理同样启用了 Mipmap 与各向异性过滤，以缓解远处与掠射角下的纹理走样——这正是实验一中超采样（SSAA）与 Mipmap 抗锯齿思想在表面细节上的延续。

% ------------------------------------------
% 第四章
% ------------------------------------------
\chapter{阴影映射与基于 FBO 的后处理特效}

\section{Shadow Mapping 阴影映射}
阴影映射由 Williams（1978）提出，是实时渲染最常用的阴影算法，分两遍渲染：第一遍以光源为视点渲染场景，仅记录每个方向最近表面的深度，得到深度图（shadow map）；第二遍以相机为视点正常渲染，将每个着色点变换到光源空间，比较其深度与深度图存储深度以判断遮挡。设光源视图矩阵为 $V_l$、投影矩阵为 $P_l$、物体模型矩阵为 $M$，则世界点 $p$ 在光源裁剪空间坐标为：
\begin{equation}
p_{\text{light}} = P_l \cdot V_l \cdot M \cdot p_{\text{model}}
\end{equation}
经透视除法与 $[0,1]$ 重映射得到采样坐标 $(u,v)$ 与当前深度 $d$，遮挡判定为：
\begin{equation}
\text{shadow} = \big(\, d - \text{bias} > \text{depthMap}(u,v) \,\big)\ ?\ 0 : 1
\end{equation}
其中 $\text{bias}$ 为深度偏移，用于消除自遮挡引起的“阴影痤疮”（shadow acne）。

\subsection{PCF 软阴影}
为缓解深度图分辨率有限导致的硬边走样，采用\keyword{百分比邻近过滤（PCF）}：在采样点邻域多次比较并取平均，得到 $[0,1]$ 的连续可见度，使阴影边缘柔和：
\begin{equation}
\text{visibility} = \frac{1}{n}\sum_{k=1}^{n} \text{compare}\big(\, d-\text{bias},\ \text{depthMap}(u+\Delta u_k,\, v+\Delta v_k)\,\big)
\end{equation}
工程上，主方向光开启阴影投射，阴影贴图分辨率设为 $2048\times2048$（高于课程要求的 $1024\times1024$），渲染器采用 \code{PCFSoftShadowMap} 软阴影类型，并通过 \code{shadow.bias} 与 \code{shadow.radius} 调校以兼顾精度与边缘平滑。参数面板提供阴影开关，可实时对比有无阴影的效果与开销。

\section{基于帧缓冲(FBO)的后处理管线}
后处理基于 Three.js 的 \code{EffectComposer} 实现，将场景渲染至离屏帧缓冲（FBO），再串联多个全屏 Pass 进行图像级处理。本项目实现了\keyword{Bloom 辉光}与\keyword{Sobel 边缘检测}两类特效，满足“至少两种后处理特效”的要求，且均支持参数面板实时调节。

\subsection{Bloom 辉光}
Bloom 模拟强光的泛光晕染：先以亮度阈值提取高亮区域，再做多级高斯模糊，最后叠加回原图。高斯模糊核为：
\begin{equation}
G(x,y) = \frac{1}{2\pi\sigma^2}\, e^{-\frac{x^2+y^2}{2\sigma^2}}
\end{equation}
本项目使用 \code{UnrealBloomPass}，可调节强度、半径与阈值。

\subsection{Sobel 边缘检测}
边缘检测以自定义 GLSL 着色器实现：在亮度图上用 Sobel 算子分别求水平与垂直梯度，合成梯度幅值作为边缘强度，叠加为轮廓线。Sobel 卷积核与梯度幅值为：
\begin{equation}
G_x = \begin{bmatrix} -1 & 0 & +1 \\ -2 & 0 & +2 \\ -1 & 0 & +1 \end{bmatrix} * L,\quad
G_y = \begin{bmatrix} +1 & +2 & +1 \\ 0 & 0 & 0 \\ -1 & -2 & -1 \end{bmatrix} * L,\quad
E = \sqrt{G_x^2 + G_y^2}
\end{equation}
其中 $L$ 为像素亮度。核心片元着色器如代码清单~\ref{lst:sobel}。

\begin{lstlisting}[language=C, caption={Sobel 边缘检测（GLSL 片元着色器核心）}, label={lst:sobel}]
float lum(vec2 uv){ vec3 c=texture2D(tDiffuse,uv).rgb;
    return dot(c, vec3(0.299,0.587,0.114)); }
// 3x3 邻域亮度做 Sobel 卷积
float gx = -tl - 2.0*l - bl + tr + 2.0*r + br;
float gy =  tl + 2.0*t + tr - bl - 2.0*b - br;
float edge = smoothstep(0.15, 0.6, sqrt(gx*gx + gy*gy));
gl_FragColor = vec4(mix(base, mix(base, uEdgeColor, edge), uStrength), 1.0);
\end{lstlisting}

\section{渲染效果展示}
基于上述管线，本项目对关键特效执行了开关对比，结果如图~\ref{fig:results} 所示。
\begin{figure}[H]
    \centering
    \begin{subfigure}[b]{0.48\textwidth}
        \centering
        \figbox{overall.jpg}
        \caption{整体场景（PBR + 阴影 + 法线贴图）}
        \label{fig:overall}
    \end{subfigure}
    \hfill
    \begin{subfigure}[b]{0.48\textwidth}
        \centering
        \figbox{blinn_pbr.jpg}
        \caption{Blinn-Phong 与 PBR 着色对比}
        \label{fig:blinn}
    \end{subfigure}

    \vspace{0.5cm}

    \begin{subfigure}[b]{0.48\textwidth}
        \centering
        \figbox{shadow.jpg}
        \caption{阴影开/关对比}
        \label{fig:shadow}
    \end{subfigure}
    \hfill
    \begin{subfigure}[b]{0.48\textwidth}
        \centering
        \figbox{postfx.jpg}
        \caption{Bloom 辉光与边缘检测后处理}
        \label{fig:postfx}
    \end{subfigure}

    \caption{室内场景在不同渲染特性下的效果对比}
    \label{fig:results}
\end{figure}
% 截图说明：请在浏览器中进入场景调好状态后按 ESC 退出锁定再截图，
% 将四张图分别命名为 overall.jpg / blinn_pbr.jpg / shadow.jpg / postfx.jpg 放在本 .tex 同目录。

\subsection{各效果图的学术解读}
对图~\ref{fig:results} 中各子图所体现的渲染原理逐一解读如下：
\begin{itemize}
    \item \keyword{图~\ref{fig:overall}（整体场景）}：在 PBR 材质、阴影映射与法线贴图全部开启下的完整画面。可见木桌、椅子在地面投下边缘柔和的 PCF 软阴影，地面与墙面在掠射光下呈现由法线贴图带来的细密凹凸，台灯点光在桌面形成局部暖色高光，整体光照层次分明。
    \item \keyword{图~\ref{fig:blinn}（Blinn-Phong 与 PBR 对比）}：左右（或切换前后）分别为自写 Blinn-Phong 着色与引擎 PBR 着色。可观察到 Blinn-Phong 因缺乏能量守恒、在强光下整体偏亮、明暗过渡较“硬”，且在桌面等平整大面上镜面高光成片出现；而 PBR 凭借微表面 BRDF 与金属/粗糙度参数，高光更聚拢、材质质感更接近真实，直观印证了经验模型与物理模型的差异。
    \item \keyword{图~\ref{fig:shadow}（阴影开/关对比）}：关闭阴影后，物体与地面失去接触感、空间关系变得扁平，同时统计面板显示 draw call 由 57 降至 12，量化体现了阴影映射“以额外渲染通路换取空间真实感”的开销-收益关系。
    \item \keyword{图~\ref{fig:postfx}（后处理特效）}：开启 Bloom 后，台灯等高亮区域产生柔和泛光晕染；开启 Sobel 边缘检测后，物体轮廓被提取为描边线，呈现非真实感（NPR）风格。二者均为基于 FBO 的离屏图像级处理，验证了后处理管线的可组合性。
\end{itemize}

% ------------------------------------------
% 第五章
% ------------------------------------------
\chapter{交互系统与实时性能量化}

\section{交互式控制}
系统实现了完整的第一人称交互：点击锁定鼠标（Pointer Lock），\code{WASD} 平移、鼠标转动视角、\code{Space}/\code{Shift} 升降，操作流畅。\keyword{射线拾取}方面，以屏幕中心准星向场景投射射线，对可拾取物体求交并高亮选中（多材质模型按逻辑物体整体高亮），同时在面板回显其材质参数。\keyword{实时参数面板}支持动态调节光照强度、阴影开关、选中物体的金属度与粗糙度，以及 Bloom 与边缘检测强度，所有修改即时生效。

\section{实时性能量化度量}
在开启不同特性的条件下对渲染器进行了帧率与 draw call 度量。为获得真实统计，关闭了 \code{renderer.info} 的自动重置并在每帧起始手动重置，以累计本帧（含后处理多 Pass 与阴影通路）的全部渲染调用。典型度量如表~\ref{tab:perf}。

\begin{table}[H]
\centering
\caption{不同渲染配置下的实时性能指标}
\label{tab:perf}
\begin{tabular}{p{4cm} p{2.5cm} p{2.5cm} p{3cm}}
\toprule
\textbf{配置条件} & \textbf{帧率(FPS)} & \textbf{Draw Calls} & \textbf{说明} \\
\midrule
基础(阴影开,无后处理) & 约 165 & 57 & 含阴影通路 \\
关闭阴影 & 约 165 & 12 & 省去阴影贴图渲染 \\
开启 Bloom & 约 165 & 70 & Bloom 多 Pass \\
开启 Blinn-Phong 对照 & 约 165 & 57 & 着色器切换无额外 Pass \\
\bottomrule
\end{tabular}
\end{table}
由表可见，阴影通路是主要的 draw call 来源（关闭后由 57 降至 12），后处理 Pass 会按特效数量线性增加 draw call，而自写 Blinn-Phong 仅替换材质、不引入额外 Pass，开销与 PBR 相当。

\section{工程难点与解决}
\begin{longtable}{|p{0.27\textwidth}|p{0.35\textwidth}|p{0.28\textwidth}|}
\hline
\textbf{遭遇障碍} & \textbf{原因剖析} & \textbf{解决对策} \\
\hline
\endhead
Blinn-Phong 在平整大面上高光“洗白”。 & 平面法向量恒定，半程向量近似不变，镜面高光整片覆盖；近距离点光叠加进一步提亮。 & 提高高光指数、降低镜面系数与点光权重，使高光收敛为紧凑亮斑；并保留该现象作为与 PBR 的对比素材。 \\
\hline
后处理导致 draw call 统计失真。 & \code{EffectComposer} 输出 Pass 覆盖 \code{renderer.info}，仅反映最后一个全屏四边形。 & 关闭 info 自动重置，每帧起始手动重置，累计本帧全部渲染调用。 \\
\hline
Blender(Z-up) 与 Three.js(Y-up) 坐标差异、多材质拆分。 & glTF 导出朝向不一致；合并体的多材质被拆为多个子网格导致高亮串色。 & 导出启用 Y-up；加载后按物体名收集网格并为每个网格克隆独立材质。 \\
\hline
\end{longtable}

% ------------------------------------------
% 第六章
% ------------------------------------------
\chapter{总结与展望}
本项目从数学原理出发，白盒化实现了 Blinn-Phong 光照、切线空间法线贴图与 Shadow Mapping 等关键实时渲染技术，并以 Blender 程序化建模管线提供全部几何资产，辅以第一人称交互、投影切换、射线拾取与基于 FBO 的后处理特效，构建了一个完整、可交互、可对比的室内高级渲染系统。本项目是对课程两次实验的综合与延伸：实验一的光栅化管线（坐标变换、纹理映射、抗锯齿）为本项目的实时渲染奠定基础，实验二的光线追踪与蒙特卡洛全局光照（BRDF、可见性、重要性采样）则为本项目的光照与阴影近似提供了理论参照。通过本次实践，本人深入理解了实时渲染管线、局部光照模型与阴影/法线技术的数学原理与工程取舍，并体会到实时光栅化与离线路径追踪在“效果—性能”权衡上的互补关系。

后续工作可从三方面深化：其一，将经验性 Blinn-Phong 升级为完整的 Cook-Torrance 微表面 BRDF，并引入基于图像的环境光照（IBL）与屏幕空间环境光遮蔽（SSAO）以增强空间感；其二，加入 LOD 细节层次与视锥/背面剔除，提升大规模场景下的帧率；其三，引入几何着色器驱动的动态网格或实例化渲染，表现植被、人群等复杂场景，进一步逼近工业级实时渲染水平。

% ------------------------------------------
% 第七章：参考文献与代码网址
% ------------------------------------------
\chapter{参考文献与开源代码网址}

\section{学术参考文献 (References)}
\begin{enumerate}
    \item[1] PHONG B T. Illumination for computer generated pictures[J]. Communications of the ACM, 1975, 18(6): 311-317.
    \item[2] BLINN J F. Models of light reflection for computer synthesized pictures[C]//Proceedings of the 4th annual conference on Computer graphics and interactive techniques (SIGGRAPH). 1977: 192-198.
    \item[3] WILLIAMS L. Casting curved shadows on curved surfaces[C]//Proceedings of SIGGRAPH. 1978: 270-274.
    \item[4] REEVES W T, SALESIN D H, COOK R L. Rendering antialiased shadows with depth maps[C]//Proceedings of SIGGRAPH. 1987: 283-291.
    \item[5] PHARR M, JAKOB W, HUMPHREYS G. Physically Based Rendering: From Theory to Implementation[M]. 4th ed. Cambridge: MIT Press, 2023.
    \item[6] AKENINE-MÖLLER T, HAINES E, HOFFMAN N, et al. Real-Time Rendering[M]. 4th ed. Boca Raton: CRC Press, 2018.
    \item[7] KAJIYA J T. The rendering equation[C]//Proceedings of SIGGRAPH. 1986: 143-150.
    \item[8] MÖLLER T, TRUMBORE B. Fast, minimum storage ray-triangle intersection[J]. Journal of Graphics Tools, 1997, 2(1): 21-28.
    \item[9] 张小红. 计算机图形学与高级渲染课程实验讲义（实验一：光栅化管线；实验二：光线追踪与全局光照）[Z]. 2025.
\end{enumerate}

\section{开源代码仓储与在线演示网址 (Code Repository \& Online Demo)}
\begin{itemize}
    \item \keyword{项目代码仓储 (GitHub Repository)}: \\
    \url{https://github.com/your-username/graphics-final} \\
    \textit{（注：包含 index.html、模块化 src/ 源码、自定义 GLSL 着色器与 Blender 建模脚本。）}
    \item \keyword{Three.js 官方文档}: \\
    \url{https://threejs.org/docs/}
\end{itemize}

\end{document}
